\section{Initial-State Implementation and Approximate Responses}
\label{sec:r19-enforcement}
Fix the benefit parameter and write $\mathcal B_r(p)=B_r(p)+dA_r(p)$ for fee-free private payoff. Let $H_r(p)\geq0$ be discounted elective surrender and let
\begin{equation}
 W_{\infty,r}=\sup_{p:H_r(p)=0}\mathcal B_r(p),\qquad
 F^{\rm init}_r=\max\left\{0,\sup_{p:H_r(p)>0}\frac{\mathcal B_r(p)-W_{\infty,r}}{H_r(p)}\right\}.
 \label{eq:r19-initial-threshold}
\end{equation}
A supremum over an empty surrender class is understood as $-\infty$. Assume that the no-surrender class is nonempty, its value is finite and attained, and $F^{\rm init}_r<\infty$. The threshold is a property of the initial contract, including its initial state and admissible mandate. It is not obtained by maximizing a gap over states that the contract cannot reach.

\begin{theorem}[Implementation and the global response-loss frontier]
\label{thm:r19-enforcement}
Under these assumptions, $W_r(F)=\sup_p\{\mathcal B_r(p)-FH_r(p)\}$ equals $W_{\infty,r}$ if and only if $F\geq F^{\rm init}_r$. If $F>F^{\rm init}_r$, every exact best response has zero elective surrender, and every global $\eta$-response obeys
\begin{equation}
 H_r(p)\leq \min\left\{H_{\max},\frac{\eta}{F-F^{\rm init}_r}\right\},
 \label{eq:r19-eta-frontier}
\end{equation}
where $H_{\max}$ is any valid upper bound on surrender. The second bound is sharp over finite dynamic economies allowing public randomization. Thus $F\geq F^{\rm init}_r+\eta/h$, with $h>0$, is sufficient to guarantee $H_r\leq h$. A claim of zero surrender for every positive-$\eta$ response generally requires more than an exact-policy exclusion test.
\end{theorem}
\begin{proof}
For $H_r(p)>0$, the definition gives $\mathcal B_r(p)\leq W_{\infty,r}+F^{\rm init}_rH_r(p)$. The same inequality holds for $H_r(p)=0$. Consequently
\begin{equation}
 \mathcal B_r(p)-FH_r(p)\leq W_{\infty,r}-(F-F^{\rm init}_r)H_r(p).
 \label{eq:r19-fee-gap}
\end{equation}
When $F\geq F^{\rm init}_r$, a no-surrender optimum attains the right value and no policy exceeds it. If $0\leq F<F^{\rm init}_r$, some positive-surrender policy has ratio greater than $F$ and strictly exceeds $W_{\infty,r}$. Strict inequality $F>F^{\rm init}_r$ in \eqref{eq:r19-fee-gap} rules out positive surrender at an exact optimum. Combining the same inequality with private payoff at least $W_{\infty,r}-\eta$ proves \eqref{eq:r19-eta-frontier}.

For sharpness, take a no-surrender action with payoff zero and a surrender action with $H=H_{\max}>0$ and fee-free payoff $F^{\rm init}H_{\max}$. A public lottery choosing the latter with probability $a$ loses exactly $aH_{\max}(F-F^{\rm init})$ and produces surrender $aH_{\max}$. Choosing the largest $a\leq1$ allowed by the loss budget attains the bound. Such a menu is a one-date dynamic economy. At the threshold the two actions tie, which also explains why weak implementation does not imply zero surrender for every best response at equality.
\end{proof}

\subsection{When adjustment ranks initial thresholds}
The no-surrender Bellman value at date and state $(n,x)$ is denoted $V_{\infty,r,n}(x)$. With the same immediate surrender payoff $G_n(x)-F$ and fee incidence in two regimes, nested operating controls imply $V_{\infty,1,n}\geq V_{\infty,0,n}$. Hence the uniform obstacle thresholds
\begin{equation}
 F^{\rm unif}_r=\sup_{(n,x)\text{ eligible}}[G_n(x)-V_{\infty,r,n}(x)]_+
 \label{eq:r19-uniform-threshold}
\end{equation}
are weakly ordered. Eligibility excludes forced discharge and includes precisely the dates and live states at which elective surrender is available. This standard verification ordering does not itself order \eqref{eq:r19-initial-threshold}. The following gives two explicit additional conditions.

\begin{proposition}[Two sufficient conditions for an initial-state ordering]
\label{prop:r19-initial-order}
\begin{enumerate}
\item Suppose that in each finite-state regime there exists an attained optimal no-surrender policy from the initial contract that visits every eligible state-date pair with positive probability. Then $F^{\rm init}_r=F^{\rm unif}_r$. Nested operating controls and unchanged surrender payoffs therefore imply $F^{\rm init}_1\leq F^{\rm init}_0$.
\item Alternatively, suppose there exist $\Delta\in\mathbb R$, $\kappa\geq0$, and a map $T$ from every positive-surrender policy in regime 1 to a policy in regime 0 such that
\begin{align}
 H_0(Tp)&=H_1(p),\qquad W_{\infty,1}-W_{\infty,0}\geq\Delta,\\
 \mathcal B_1(p)-\mathcal B_0(Tp)&\leq\Delta-\kappa H_1(p).
 \label{eq:r19-gain-dominance}
\end{align}
Then $F^{\rm init}_1\leq\max\{0,F^{\rm init}_0-\kappa\}$.
\end{enumerate}
Neither condition is asserted merely because one action set includes the other.
\end{proposition}
\begin{proof}
For the first statement, $F\geq F^{\rm unif}_r$ makes the no-surrender Bellman value dominate every elective stopping obstacle. Backward induction makes it also the value with stopping, so $F^{\rm init}_r\leq F^{\rm unif}_r$. If $0\leq F<F^{\rm unif}_r$, some eligible pair has a strictly profitable stopping gap. Follow the hypothesized no-surrender optimum until that pair and stop there. The pair is reached with positive probability, and optimal continuation there equals $V_{\infty,r,n}(x)$ because the policy is optimal at every reached history of positive probability. The modification strictly increases initial payoff. Thus $F<F^{\rm init}_r$; taking the supremum proves equality. The same argument extends to continuous states when continuity of the gap and positive occupation of every relevant open neighborhood replace positive mass at a point.

For the second statement, subtract the inequality for no-surrender values from \eqref{eq:r19-gain-dominance}, divide by the common positive $H$, and obtain
\[
 \frac{\mathcal B_1(p)-W_{\infty,1}}{H_1(p)}
 \leq \frac{\mathcal B_0(Tp)-W_{\infty,0}}{H_0(Tp)}-\kappa
 \leq F^{\rm init}_0-\kappa.
\]
Take the supremum and the positive part. This condition compares the adjustment gain on surrender deviations with the gain from continuing; it is stronger than a comparison of no-surrender values alone.
\end{proof}

\subsection{Value-message implementation without a selected stopping rule}
For $\delta>0$ and any $\eta$-response at fee $F$,
\begin{equation}
 H_r(p)\leq\frac{W_r(F-\delta)-W_r(F)+\eta}{\delta}
 \leq\frac{U_r(F-\delta)-L_r(F)+\eta}{\delta}.
 \label{eq:r19-query-exit}
\end{equation}
The first inequality follows by evaluating the same policy at the lower fee. It remains valid for global near-optimal behavior and charges both value errors before reporting surrender. When $W_r(F)=W_{\infty,r}$ has separately been established, the corresponding no-surrender lower bound can be used for $L_r(F)$. A fee grid cannot certify an exact zero derivative simply because two floating-point values coincide. Our numerical implementation instead verifies obstacle inequalities on all states and a strictly feasible surrender witness below the reported lower threshold. It reports a bracket, its arithmetic allowance, the first-date attainment convention, and the type of threshold certified.
